\documentclass[../main.tex]{subfiles}
\graphicspath{{\subfix{../../images/}}}

\begin{document}

\section{Usage}

Once the \cc{notex} CLI utility is installed, it offers a set of commands to facilitate managing and working with your \LaTeX{} projects.
The following table summarizes the main commands and their purpose. Further details can be obtained by running \cc{notex help} (or \cc{notex -h}).

\vspace{4mm}
\begin{center}
\renewcommand{\arraystretch}{1.4}
\begin{tabular}{lp{0.53\textwidth}}
    \toprule
    \textbf{Command} & \textbf{Description} \\
    \midrule
    \cc{notex help}             & Shows \cc{notex} binary options and help guide. \\
    \cc{notex version}          & Shows current binary version. \\
    \cc{notex info}             & Displays installation and current project information. \\
    \cc{notex init .}           & Initializes a project in the current working directory. \\
    \cc{notex clean <path>}     & Cleans recursively temporary files from the provided \cc{<path>}. \\
    \cc{notex get <snippet>}    & Provides a ready-to-use \LaTeX{} snippet template. \\
    \cc{notex theme <theme>}    & Changes the current project theme. \\
    \cc{notex add <what>}       & Adds a new section or bibliography to the current project. \\
    \cc{notex remove <what>}    & Removes a section or bibliography from the current project. \\
    \cc{notex ls}               & Lists the files that belong to the current project. \\
    \cc{notex checkhealth}      & Checks the \TeX{} environment, installation, and project state. \\
    \cc{notex set <key> <what>} & Changes a whitelisted \cc{notex.json} key directly. \\
    \cc{notex reset}            & Regenerates \cc{main.tex} from the template. \\
    \cc{notex delete}           & Removes notex files and build artefacts, keeps your documents. \\
    \bottomrule
\end{tabular}
\end{center}
\vspace{2mm}

\begin{info}
    This guide assumes that the \cc{notex} utility has been installed as a globally available binary.
\end{info}

\subsection{Initialization}

You can quickly set up a \LaTeX{} project in the current directory (\itt{or in one of its sub-paths}) via the \bft{init command}:

\begin{bash}
    notex init .
    # or equivalent:
    # notex init multi .
\end{bash}

The default initialization starts a \bft{multi-file} project: a \cc{main.tex} using the \cc{subfiles} package is created together with a \cc{./sections/} folder for individual section files.
An example of this layout can be found in the \cc{examples/multi/} directory.

For smaller projects, a \bft{mono-file} project with a single \cc{main.tex} can be initialized instead:

\begin{bash}
    notex init mono .
\end{bash}

An example can be found in the \cc{examples/mono/} directory.

Both commands also create a \cc{.notex/} directory at the project root containing NoTeX configuration.

\begin{info}
    \cc{notex init} refuses to overwrite an existing \cc{main.tex} or \cc{.notex/} directory.
    Pass \cc{--force} to overwrite anyway.
\end{info}

In both cases, the generated \cc{main.tex} automatically uses\\
 \cc{\textbackslash documentclass\{settings/notex\}} when a local installation is detected in the target directory, and \cc{\textbackslash documentclass\{notex\}} otherwise.

\subsection{Management}

Once you have a working project, \cc{notex} provides a series of management commands to facilitate common operations.

\subsubsection{Change Theme}

You can change the theme used by the \LaTeX{} template with:

\begin{bash}
    notex theme <theme_name>
\end{bash}

where \cc{<theme\_name>} can be one of the available themes:

\begin{itemize}
    \item \cc{light} (\itt{default})
    \item \cc{dark}
    \item \cc{tokyo}
    \item \cc{bw}
\end{itemize}

\subsubsection{Manage Sections and Bibliography}

You can add a new section with:

\begin{bash}
    notex add section "My Title"
\end{bash}

This adds a section directly to \cc{main.tex} for a mono-file project, or creates a new section file under \cc{sections/} for a multi-file project.

You can remove a section by its number \cc{<N>}:

\begin{bash}
    notex remove section <N>
\end{bash}

Two analogous commands exist for bibliography management.\\
To create \cc{bibliography.bib} and wire it into \cc{main.tex}:

\begin{bash}
    notex add bib
\end{bash}

To remove the bibliography:

\begin{bash}
    notex remove bib
\end{bash}

\subsubsection{Cleaning}

\LaTeX{} projects accumulate temporary files and build artefacts. \cc{notex} provides a command to clean them up recursively:

\begin{bash}
    notex clean <path>
\end{bash}

\begin{warning}
    Cleaned files are \bft{lost forever} once the command runs. To preview which files would be removed \bft{before actually deleting them}, run a \itt{dry run}:

    \begin{bash}
        notex clean --dry-run
    \end{bash}

    This lists the files that would be removed by an actual clean operation.
\end{warning}

\subsubsection{References and Snippets}

To recall the \LaTeX{} environments provided by the NoTeX template, the \cc{get} command prints ready-to-paste snippet templates:

\begin{bash}
    notex get
\end{bash}

Lists every available snippet name (\itt{info, warning, cbox, ebox, code, \ldots}).
To fetch a specific snippet, pass its name:

\begin{bash}
    notex get box       # displays a template for the "box" environment
    notex get warning   # displays a template for the warning callout
    notex get code      # displays a template for the code environment
\end{bash}

\subsection{State and Info}

Several commands allow inspecting and modifying the current project state.

\begin{bash}
    notex ls
\end{bash}

Lists the files that belong to the project (\itt{main file, sections, bibliography}).

\begin{bash}
    notex checkhealth
\end{bash}

Checks the \TeX{} environment, the installation, and the project's state.

\begin{bash}
    notex set main_file alt.tex
\end{bash}

Changes a whitelisted \cc{notex.json} key directly (\itt{e.g.} \cc{main\_file}, \cc{theme},\\ \cc{bibliography\_file}).

\begin{bash}
    notex reset
\end{bash}

Regenerates \cc{main.tex} from the template, backing up the existing file to \cc{main.tex.bak}.

\begin{bash}
    notex delete
\end{bash}

Removes \cc{.notex/}, \cc{settings/}, \cc{fonts/}, and build artefacts while keeping your documents untouched.
To remove the entire project scaffolding destructively:

\begin{bash}
    notex delete --all
\end{bash}

\end{document}


